\documentclass{article}
\usepackage[utf8]{inputenc}
\usepackage[T1]{fontenc}
\usepackage[margin=2cm]{geometry}
\usepackage{hyperref}
\sloppy
\begin{document}
\section*{Historical source}
Encyclopaedia Britannica scan, source row 2424. Parse 3.5 extracted the text. Super 3.5 VL corrected it from an enlarged crop of each block. It can still contain recognition errors.
Page headers, footers, and page numbers are removed.

\subsection*{AGRICULTURE (E0216)}
\textit{Continues an entry from an earlier page.}

Live Stock. The lid is then closed for a short time, and the mess thereafter served out to the cattle. The meal may be that of barley, oats, beans, or, in short, whatever is cheapest at the time, including all the inferior grain of the farm, with the addition of ¼ lb. of ground linseed for each animal. An appropriate mode of administering the food is to give 50 lb. of sliced turnips early in the morning, the cooked mess about noon, and as much more turnips in the afternoon. A little oat straw is also given to them daily. For a few weeks before the completion of the fattening process, it is usually remunerative to give 4 or 5 lb. of linseed cake, and 6 lb. of good hay to each beast daily; the cake in combination with the cooked mess, and the hay in lieu of the oat-straw. From the peculiar structure of the digestive organs of the ox—which are adapted for disposing of bulky and but moderately nutritious food—this plan of inducing him to eat a considerable quantity of straw promotes his health and comfort, and economises the more costly food. The Rev. Mr Huxtable has pointed out a still cheaper way of accomplishing the same end. He shreds his turnips by a machine, analogous to a bark-mill, and then mixes them with cut straw, so that the cattle must eat both together. His opinion is, that he not only gains the element of bulk in this way, but that, from the straw and turnips being swallowed together, the hurtful chilling effect of the latter is avoided. He finds that 70 lb. of turnips thus administered is a sufficient daily ration for a fattening bullock.

Best mode of housing cattle.

A better appreciation of the effects of temperature on the animal economy has of late years exerted a beneficial influence upon the treatment of fattening cattle. Observant

\textbackslash{}begin\{tabular\}\{cccc\}
best mods\&of housing\&stable." Sheep is the most\textbackslash{}\textbackslash{}
attablet. combats\&costs have not two\&since they are expressed an opinion in which agriculturists of the present day fully con- can. But if this was we use of the flocks of his sheep, how much\&more of the many animals breeds which now are the rich features, the grassy downs, and healthful mountains of our native countries. Their birds is high estimated with\textbackslash{}\textbackslash{}
\textbackslash{}end\{tabular\}

1st, Facilities for supplying food and litter, and for removing dung with the utmost economy of time and labour.

2d. Complete freedom from disturbance.

2d. Complete freedom from disturbance.
3d. A moderate and equal degree of warmth.
4th. A constant supply of pure air.

4th, A constant supply of pure air.

4th, A constant supply of pure air.
5th, Opportunity for the cattle having a slight degree of
exercise; and
6th, The production of manure of the best quality.

6th, The production of manure of the best quality.

The production of manure of the best quality. We have no hesitation in expressing our opinion that the whole of these conditions are attained most fully by means of well-arranged boxes. Stalls are to be preferred where the saving of litter is an object; and yards for the rearing of young cattle, which require more exercise than is suitable for fattening stock.

for fattening stock.
But however excellent the system of housing and feeding
which is adopted, a successful result will, in every instance,
be much dependent upon the vigilant superintendence of the

owner, and his skill in so managing the commissariat, as to secure throughout the year a sufficiency of suitable food for the stock on hand. Unless this is attended to, he may find himself necessitated to sell his cattle before they are fat, and when markets are glutted. Whenever they have attained to what is technically called ripe fatness, and prices are at their average rate, it is generally more profitable forthwith to dispose of them, than to speculate upon a rise in the markets.

\subsection*{SIEER (E0217)}
SIEER.

When Fitzherbert so long ago said, "Sheep is the most profitable cattle that a man can have," he expressed an opinion in which agriculturists of the present day fully concur. But if this was true of the flocks of his time, how much more of the many admirable breeds which now adorn the rich pastures, the grassy downs, and heath-clad mountains of our native country. Their flesh is in high estimation with all classes of the community, and constitutes at least one-half of all the butcher meat consumed by them. Their fleeces supply the raw material for one of our most flourishing manufactures. They furnish to the farmer an important source of revenue, and the readiest means of maintaining the fertility of his fields. The distinct breeds, and sub-varieties of sheep found in Great Britain are very numerous. We have no intention of noticing them in detail, but shall rather confine our observations to those which, by common consent, are the most valuable for their respective and appropriate habitats. They may be fitly classed under these three heads, viz. the heavy breeds of the low country, those found on our downs and similar localities, and the mountain breeds.

Of the first class, the improved Leicesters are still the most important to the country. They are more widely diffused in the kingdom than any of their congeners. Although, from the altered taste of the community, their mutton is less esteemed than formerly, they still constitute the staple breed of the midland counties of England. Leicester rams are also more in demand than ever, for crossing with other breeds. It is now about a century since this breed was produced by the genius and perseverance of Bakewell, in whose hands they attained a degree of excellence that has probably not yet been exceeded by the many who have cultivated them since his day. The characteristics of this breed are extreme docility, extraordinary aptitude to fatten, and the early age at which they come to maturity. The most marked feature in their structure is the smallness of their heads, and of their bones generally, as contrasted with their weight of carcase. They are clean in the jaws, with a full eye, thin ears, and placid countenance. Their backs are straight, broad, and flat, the ribs arched, the belly carried very light, so that they present nearly as straight a line below as above, the chest is wide, the skin very mellow, and covered with a beautiful fleece of long, soft wool, which weighs on the average from 6 to 7 lb. On good soils, and under careful treatment these sheep are currently brought to weigh from 18 to 20 lb. per quarter at 14 months old, at which age they are now usually slaughtered. At this age their flesh is tender and juicy; but when carried on until they are older and heavier, fat accumulates so unduly in proportion to the lean meat, as to detract from its palatableness and market value.

Lincolns.—These were at one time very large, ungainly Lincoln animals, with an immense fleece of very long wool. By breeding, crossing them with the Leicesters the character of the breed has been entirely changed, and very greatly for the better. It is now, in fact, a subvariety of the Leicester, with larger frame and heavier fleece than the pure breed. Sheep of

See "Present Prices" by Rev. A. Haxtable. Ridgeway. London, 1850.


\section*{Source record}
\begin{description}
\item[Scan ID] \texttt{britannica:309f765b1ecc83e0a7f7762e44b4a927a7555443:train:2424}
\item[Image file] \texttt{page-2424-dd2506be16cf.jpg}
\item[Image SHA-256] \texttt{dd2506be16cf7abc83d6915d2eeca90ede40d89a0e648c6741d1ea42ee73e4b7}
\item[Dataset] NationalLibraryOfScotland/encyclopaedia\_britannica\_illustrated at revision \texttt{309f765b1ecc83e0a7f7762e44b4a927a7555443} (CC0-1.0)
\item[Extraction] model \texttt{nvidia/nemotron-3.5-parse}, raw reply \texttt{extract/raw/page-2424-dd2506be16cf.json}
\item[Raw reply SHA-256] \texttt{f422f1a5774c89c7d9635357189295be73557d1e705a331d2527ce7667fe2cb6}
\item[Correction] model \texttt{nvidia/Nemotron-3.5-Super-VL}, block crops, raw replies \texttt{correct/raw/crop\_nothink/page-2424-dd2506be16cf-b*.json}; parse: 1, parse\_fallback: 1, super\_vl: 15
\item[Running header] AGRICULTURE. 245
\end{description}
\end{document}
